\documentclass{scrartcl}
\usepackage[T1]{fontenc}
\usepackage[ngerman]{babel}
\usepackage{graphicx}
\usepackage[hidelinks]{hyperref}
\usepackage{subcaption} %erforderlich für subfigure

\usepackage{tikz}
\usepackage{float}
\usetikzlibrary{calc}

\usetikzlibrary{intersections}
\usetikzlibrary{arrows}

\usetikzlibrary{
	calc,
	intersections,
	arrows.meta,
	shapes,
	quotes,
	angles,
	patterns
}


\begin{document}

	% =========================================
	\section{Parametrische Zeichnungen mit TikZ}
	% =========================================
	

%\usetikzlibrary{calc,patterns,angles,decorations.markings,babel,quotes}


\selectlanguage{english}

\begin{figure}[htbp]
	\centering

\begin{tikzpicture}

\coordinate (origo) at (0,0);
\coordinate (pivot) at (1,5);

% draw axes
\fill[black] (origo) circle (0.1);
\draw[thick,gray,->] (origo) -- ++(4,0) node[black,right] {$x$};
\draw[dashed,thick,gray,-] (origo) -- ++(0,-4);% node (mary) [black,below] {$y_1$};

% draw roof
\fill[pattern = north east lines] ($ (origo) + (-1,0) $) rectangle ($ (origo) + (1,0.5) $);
\draw[thick] ($ (origo) + (-1,0) $) -- ($ (origo) + (1,0) $);

\draw[thick] (origo) -- ++(300:3) coordinate (bob);
\fill (bob) circle (0.2) node[right=0.2cm] {$m_1$};

%\draw[thick] (origo) -- ++(300:3) coordinate (bob);
\draw[thick] (300:3) -- ++(330:2) coordinate (bob1);
\fill (bob1) circle (0.2) node[right=0.2cm] {$m_2$};

%\draw pic["$\theta$", draw, angle radius=1.2cm] {angle = mary--origo--bob};
%\draw pic["$\theta$",draw=red,-,angle eccentricity=1.50,angle radius=1cm] {angle = mary--origo--bob};

\draw [<->] (-90:1.5)  arc (-90:-60:1.5) node [anchor=north,pos=0.5] {$\theta_1$};

\draw[dashed,thick,gray,-] (300:3) coordinate (bob) -- ++(0,-2);% node (mary) [black,below] {$y_2$};
\draw [<->]  ([shift=(-90:1.5)]bob)  arc (-90:-30:1.5) node [anchor=north,pos=0.5] {$\theta_2$};
%([shift=(t:r)] x, y) is the proper starting point, where (x,y) is the center and (t:r) is the polar coordinate of starting point.
%\draw [<->] (-30:1.5)  arc (-90:-60:1.5) node [right,pos=0.5] {$\theta$};
%\draw (axis cs:.125,0)arc[radius=.25cm,start angle=0,end angle=45];

\end{tikzpicture}


    \caption{Doppelpendel mit Winkeln und Massen}
\label{fig:tikzbeispiel}

\end{figure}

Wie in Abbildung~\ref{fig:tikzbeispiel} dargestellt,
besteht das Doppelpendel aus zwei Massen.	

\end{document}
